\documentclass[10pt,a4paper]{amsart}
\usepackage[margin=19mm]{geometry}
\usepackage{amsmath,amssymb,mathtools}
\usepackage[scr=rsfs,cal=euler]{mathalfa}
\usepackage{xfrac}
\usepackage[unicode,hidelinks,bookmarks=false]{hyperref}
\newcommand{\ie}{i.e.\ }
\newcommand{\cf}{cf.\ }
\newcommand{\resp}{respectively\ }
\newcommand{\Subsection}{\subsection}
\providecommand{\nobreakdash}{\nobreak\hbox{-}}
\setlength{\emergencystretch}{2em}
\multlinegap0pt
\usepackage{amssymb}
\newtheorem{lemma}{Lemma}
\newcommand{\RR}{{\mathbb{R}}}
\newcommand{\ZZ}{{\mathbb{Z}}}
\title{On the Equivalence of the Graph-Structural and Optimization-Based Characterizations of Popular Matchings}
\author{Yuga Kanaya, Kenjiro Takazawa}
\date{Source: arXiv:2508.00349v2 (CC BY 4.0)}
\begin{document}
\maketitle

\noindent\emph{Source and licence.} On the Equivalence of the Graph-Structural and Optimization-Based Characterizations of Popular Matchings. Version arXiv:2508.00349v2 (CC BY 4.0), distributed by arXiv under the Creative Commons Attribution 4.0 International licence, http://creativecommons.org/licenses/by/4.0/, as recorded in the arXiv licence metadata for this exact version and shown on the abstract page https://arxiv.org/abs/2508.00349v2. Full licence notice: \texttt{licenses/arxiv-2508.00349-CC-BY-4.0.txt}.\par\smallskip

\noindent\textit{Editorial scope.} Counted unit: the article's lemma LEMy01 (source0.tex lines 975--979). The article's complete original proof of it is delivered with it (source0.tex lines 982--1052, one \texttt{proof} environment of 71 lines). No mathematics is written, repaired or re-derived: the statement and the proof are the article's own lines, in the article's own notation, and no retained line is substituted or re-flowed.  Retained uncounted as the source-local context the proof needs: lines 769--789; lines 959--967.  Nothing was omitted from any retained span, so the package carries no omission record.  No retained line contains a citation, so the wrapper carries no bibliography and no bibliography entry.  No retained line contains a figure environment, an image-inclusion command or drawing code, so no image is delivered and the package declares no asset dependency.  Editorial formatting only, in addition: an AMS \texttt{amsart} preamble that restates the article's own declarations and macros used by the retained lines, namely the environments \texttt{lemma} on separate counters, together with the standard packages those lines need.  The retained environments print in this document as Lemma 1 in order of appearance, whereas the article prints the same items with the numbering of its own full document.  The complete native source of the article as distributed in the arXiv e-print for this exact version is bundled unchanged as \texttt{sources/182539/source0.tex}.
The following linear program with variable $x\in \RR^E$ is a linear relaxation of the 
maximum-weight $A$-perfect matching problem in the weighted graph $(G,w_M)$: 
\begin{alignat}{3}
\label{EQLPobj}
&\mbox{Maximize}    &\quad& \sum_{(e)\in E} w_M(e)\cdot x(e)     && \\
&\mbox{subject to}  &\quad& \sum_{h\in \Gamma(a)}x(a,h)=1                &\quad& (a\in A),
\label{EQLPconst1} \\
&                   &\quad& \sum_{a\in \Gamma(h)}x(a,h)\le 1             &\quad& (h\in H),
\label{EQLPconst2}\\
&                   &\quad&x(e)\ge 0                                        &\quad& (e\in E).
\label{EQLPconst3}
\end{alignat}
The following linear program with variable $y\in \RR^{A\cup H} $ is the dual problem of \eqref{EQLPobj}--\eqref{EQLPconst3}: 
\begin{alignat}{3}
\label{EQdualobj}
&\mbox{Minimize}    &\quad& \sum_{a\in A} y(a) + \sum_{h\in H} y(h)     && \\
\label{EQdualconst1}
&\mbox{subject to}  &\quad& y(a) + y(h) \ge w_M(e)                    &\quad& (e=(a,h)\in E),\\
\label{EQdualconst2}
&                   &\quad& y(h)\ge 0                                   &\quad& (h\in H).
\end{alignat}

Note that $\chi_M(e) > 0$ implies $e\in M$, and hence $w_M(e)=1$. 
Now 
the complementary slackness condition of these linear programs is that 
\begin{alignat}{2}
    \label{EQCS1}
    &{}\mbox{$\chi_M(e) > 0$ implies $y^*(a)+y^*(h)= w_M(e)=1$}& \quad &(e=(a,h)\in E),\\
    \label{EQCS2}
    &{}\mbox{$y^*(h) > 0$ implies $\displaystyle\sum_{a\in \Gamma(h)}\chi_M(a,h)= 1$}& \quad& (h\in H).
\end{alignat}

\begin{lemma}
\label{LEMy01}
    For an integral optimal solution $y^* \in \ZZ^{A\cup H}$ of \eqref{EQdualobj}--\eqref{EQdualconst2}, 
    it holds that $y^* \in \{0,1\}^{A\cup H}$. 
\end{lemma}

\begin{proof}
    We first show that $y^*(a)\in \{0,1\}$ for each applicant $a\in A$. 
    Let $a\in A$ and $h=M(a)$. 
    It follows from $(a,h)\in M$ and \eqref{EQCS1} that $y^*(a)+y^*(h)=1$. 
    From \eqref{EQdualconst2}, we obtain that 
    \begin{align}
        \label{EQyale1}
        y^*(a)\leq 1. 
    \end{align} 

    Suppose that $h\neq l(a)$. 
    We have that $(a,l(a)) \not\in M$ and $l(a)$ is only adjacent to $a$ in $G$, 
    and hence 
    \begin{align*}
        \sum_{a'\in \Gamma(l(a))}\chi_M(a',l(a))= 0. 
    \end{align*}
    It then follows from \eqref{EQCS2} that 
    \begin{align}
    \label{EQyla}
        y^*(l(a))=0. 
    \end{align}
    We also have that 
    %%$h\in \Gamma(a) \setminus \{ l(a)\}$, 
    %%and hence 
    $h \succ_a l(a)$, 
    implying that $w_M(a,l(a))=0$. 
    It then follows from \eqref{EQdualconst1} that 
    \begin{align}
        \label{EQ3.11}
        y^*(a) + y^*(l(a)) \ge 0.
    \end{align}
    From \eqref{EQyla} and \eqref{EQ3.11}, 
    we obtain $y^*(a)\ge 0$. 
    We thus derive from \eqref{EQyale1} 
    that 
    $y^*(a)\in \{0,1\}$, 
    since $y^*$ is integral. 
    From \eqref{EQCS1}, 
    we also obtain $y^*(M(a)) \in \{0,1\}$. 

    Suppose that $h=l(a)$. 
    Let $h'\in \Gamma(a) \setminus \{l(a)\}$, 
    which must exist since $a$ is not isolated before $l(a)$ is attached. 
    We have that $h' \succ_a l(a)=M(a)$, and hence $w_M(a,h')=2$. 
    It follows from \eqref{EQdualconst1} that 
    \begin{align}
        \label{EQ3.14}
        y^*(a) + y^*(h')\ge 2, 
    \end{align}
    and further 
    \begin{align}
        \label{EQyh}
        y^*(h')\ge 1
    \end{align} from \eqref{EQyale1}. 
    It then follows from \eqref{EQCS2} that $\sum_{a\in \Gamma(h')}\chi_M(a,h')= 1$, 
    implying that $(a',h')\in M$ for some applicant $a'\in A$. 
    We can now apply the above argument to $a'$ to obtain $y^*(h')=y^*(M(a'))\in \{0,1\}$, 
    and hence $y^*(h')= 1$ from \eqref{EQyh}. 
    We thus obtain $y^*(a)= 1$ from \eqref{EQyale1} and \eqref{EQ3.14}. 

    We have shown that $y^*(a)\in \{0,1\}$ for each $a\in A$. 
    We complete the proof by showing that $y^*(h)\in \{0,1\}$ for each $h\in H$. 
    Let $h\in H$. 
    If $(a,h)\in M$ for some $a\in A$, 
    it follows from \eqref{EQCS1} that $y^*(h)=1-y^*(a)\in \{0,1\}$. 
    Otherwise, 
    $\sum_{a\in \Gamma(h)}\chi_M(a,h)= 0$, 
    and 
    we obtain $y^*(h)=0$ from \eqref{EQCS2}. 
    %%We thus conclude that $y^*(h)\in \{0,1\}$ for each $h\in H$. 
\end{proof}

\end{document}
